\section{连续时间 Score Matching}

\subsection{训练目标}

在连续时间框架下，score matching 的训练目标可以统一写为：

$$
\mathcal{L} = \mathbb{E}_{t \sim \mathcal{U}(0,T)}\;\mathbb{E}_{\mathbf{x}_0 \sim p_{\text{data}}}\;\mathbb{E}_{\mathbf{x}_t \sim q(\mathbf{x}_t|\mathbf{x}_0)}\left[\lambda(t)\left\|\mathbf{s}_\theta(\mathbf{x}_t, t) - \nabla_{\mathbf{x}_t}\log q(\mathbf{x}_t|\mathbf{x}_0)\right\|^2\right]
$$

其中：
\begin{itemize}
    \item $t$ 从 $[0, T]$ 上均匀采样（连续时间）；
    \item $\lambda(t)$ 是时间依赖的加权函数；
    \item 条件 score $\nabla_{\mathbf{x}_t}\log q(\mathbf{x}_t|\mathbf{x}_0) = -\frac{\mathbf{x}_t - \alpha_t \mathbf{x}_0}{\sigma_t^2}$ 是已知的解析形式。
\end{itemize}

\begin{importantbox}{离散 vs 连续训练}
在离散 DDPM 中，时间步 $t$ 从 $\{1, 2, \ldots, T\}$ 中均匀采样。在连续框架中，$t$ 从连续区间 $[0, T]$ 上采样。实际训练时，两者几乎没有区别——只需将离散的时间索引替换为连续的时间值即可。网络 $\mathbf{s}_\theta(\mathbf{x}, t)$ 接受连续的 $t$ 作为输入（通常通过正弦位置编码处理）。
\end{importantbox}

\subsection{加权函数 $\lambda(t)$ 的选择}

不同的 $\lambda(t)$ 选择对应不同的训练策略：

\begin{itemize}
    \item $\lambda(t) = g(t)^2$：对应 SDE 的似然加权，与变分下界（VLB）最优化相关；
    \item $\lambda(t) = \sigma_t^2$：最常见的选择，等价于噪声预测目标 $\|\boldsymbol{\epsilon}_\theta - \boldsymbol{\epsilon}\|^2$；
    \item $\lambda(t) = 1$：直接匹配 score，在低噪声时刻（小 $t$）可能不稳定。
\end{itemize}

\begin{knowledgebox}{从 Score Matching 到噪声预测}
利用 $\nabla_{\mathbf{x}_t}\log q(\mathbf{x}_t|\mathbf{x}_0) = -\boldsymbol{\epsilon}/\sigma_t$ 和 $\mathbf{s}_\theta = -\boldsymbol{\epsilon}_\theta / \sigma_t$，score matching loss 可以改写为：
$$
\left\|\mathbf{s}_\theta - \nabla_{\mathbf{x}_t}\log q\right\|^2 = \frac{1}{\sigma_t^2}\left\|\boldsymbol{\epsilon}_\theta - \boldsymbol{\epsilon}\right\|^2
$$
因此，选择 $\lambda(t) = \sigma_t^2$ 就恰好消除了 $1/\sigma_t^2$ 因子，得到简洁的噪声预测 MSE 损失——这正是 DDPM 原始论文中使用的简化目标。
\end{knowledgebox}

\subsection{本章小结}

连续时间框架下的 score matching 训练目标与离散 DDPM 的训练实质相同，只是时间变量从离散变为连续。加权函数 $\lambda(t)$ 的选择影响训练动态，但最常用的噪声预测目标在两种视角下完全一致。


